\mode*
\section{Conclusions}

\rqmisconceptions*
So far: the seed and systematic rounds document the branch-as-directory
misconception and general course-observed confusions
\autocite{IsomottonenCochez2014}, design-level concept misfits
\autocite{DeRossoJackson2016Gitless}, the commit-as-save conflation in
repository data \autocite{Cochez2013DVCSUsage}, and an emerging
multi-institution misconception catalogue
\autocite{Postner2026GitMisconceptions}; the measurement round adds Git's
hidden dependencies as professionals experience them
\autocite{Church2014HiddenDependencies} and a teacher-reported error list
\autocite{Eraslan2020GitErrors}, and finds that learners' understanding
of Git has never been measured with an instrument built for it nor
studied phenomenographically (\cref{sec:protocol-measurement}).

\rqaspects*
Twelve candidate aspects, listed per objective in
\cref{sec:model-aspects,sec:staging-aspects,sec:branches-aspects,sec:remotes-aspects}
and summarised in \cref{tab:aspects}.
Their evidential status is uneven and is stated with each: six are read
off difficulties documented in the literature, and six are posited from
the structure of the object of learning; none yet rests on our own
course data, which is still to be coded (\cref{sec:model}).
Every aspect has an item in the pre/post-test of \cref{app:quiz} (three
have open twins as well: \emph{a commit is a snapshot}, since its closed
form is asked only at the end; \emph{a branch is a pointer} and
\emph{local commits are invisible until pushed}, so that the account and
the choice can be compared), which is what will turn \enquote{candidate}
into \enquote{critical} or not.

\begin{table}
  \begin{sidecaption}{The candidate critical aspects, what each is read
    off, and the item of \cref{app:knowledge-items} that tests
    it.}[tab:aspects]
    \small
    \begin{tabular}{@{}>{\raggedright\arraybackslash}p{0.36\linewidth}l>{\raggedright\arraybackslash}p{0.34\linewidth}@{}}
      \toprule
      Aspect & Read off & Item \\
      \midrule
      \multicolumn{3}{@{}l}{\emph{The commit model} (\cref{sec:model-aspects})} \\
      A commit is a snapshot, not a diff & posited & what a commit keeps (closed, end only; open twin at both) \\
      History is a graph & literature & where earlier versions live \\
      Nothing happens without a commit & literature & after you save \\
      \addlinespace
      \multicolumn{3}{@{}l}{\emph{The staging area} (\cref{sec:staging-aspects})} \\
      Three places, two moves & literature & what git add does \\
      The same file can differ in all three places & posited & two versions at once \\
      git status reads the three places & posited & what status shows after an edit \\
      \addlinespace
      \multicolumn{3}{@{}l}{\emph{Branches} (\cref{sec:branches-aspects})} \\
      A branch is a pointer, not a place & literature & seeing another branch (closed and open) \\
      The graph outlives the pointers & posited & deleting a branch \\
      Merging joins histories & posited & what merging does \\
      \addlinespace
      \multicolumn{3}{@{}l}{\emph{Remotes} (\cref{sec:remotes-aspects})} \\
      Every clone is complete & posited & offline with a clone \\
      Push and pull move commits & literature & what git push sends \\
      Local commits are invisible until pushed & literature & why they cannot see it (closed and open) \\
      \bottomrule
    \end{tabular}
  \end{sidecaption}
\end{table}

\rqpatterns*
Preliminary patterns accompany each aspect; the branch chapter's
contrast---switching branches while \texttt{pwd} stays put---is the
paradigm case, refuting the documented misconception directly.

The natural design for validating the aspects is a learning study
\autocite[Ch.~7]{NCOL} per objective, with the pre- and post-test of
\cref{app:quiz} as its instrument: the closed items test the aspects
argued for here, and the open accounts---the first phenomenographic data
on Git, since none exists (\cref{sec:protocol-measurement})---are where
aspects we did not think of will show themselves.
